\begin{mistake}{2026-06-17}{06}

\begin{problemcontent}
设 \(y = f(x)\) 在 \(x_0\) 的某邻域内有四阶连续导数，且 \(f'(x_0) = f''(x_0) = f'''(x_0) = 0\)，且 \(f^{(4)}(x_0) < 0\)，则（\quad）

A. \(f(x)\) 在 \(x_0\) 处取得极小值\\
B. \(f(x)\) 在 \(x_0\) 处取得极大值\\
C. \((x_0, f(x_0))\) 是 \(y = f(x)\) 的拐点\\
D. \(f(x)\) 在 \(x_0\) 的某邻域内单调减少
\end{problemcontent}

\begin{solutioncontent}
正确答案：B。

\textbf{极值判断：}将 \(f(x)\) 在 \(x_0\) 处利用带拉格朗日余项的泰勒公式展开到三阶：
\[ f(x) = f(x_0) + f'(x_0)(x - x_0) + \frac{f''(x_0)}{2!}(x - x_0)^2 + \frac{f'''(x_0)}{3!}(x - x_0)^3 + \frac{f^{(4)}(\xi)}{4!}(x - x_0)^4 \]
代入导数为 \(0\) 的条件化简得：
\[ f(x) - f(x_0) = \frac{f^{(4)}(\xi)}{4!}(x - x_0)^4 \]
由于 \(f^{(4)}(x)\) 在 \(x_0\) 处连续，且 \(f^{(4)}(x_0) < 0\)，由连续函数的局部保号性可知，在 \(x_0\) 的充分小邻域内 \(f^{(4)}(\xi) < 0\)。\\
对于 \(x \neq x_0\)，\((x - x_0)^4 > 0\)，因此 \(f(x) - f(x_0) < 0\)，即 \(f(x) < f(x_0)\)。由极值定义，\(f(x)\) 在 \(x_0\) 处取得极大值。

\textbf{排错（选项 D）：}对一阶导数 \(f'(x)\) 展开：
\[ f'(x) = \frac{f^{(4)}(\eta)}{3!}(x - x_0)^3 \]
当 \(x < x_0\) 时，\((x - x_0)^3 < 0\)，又因 \(f^{(4)}(\eta) < 0\)，故 \(f'(x) > 0\)，此时单调增加；反之 \(x > x_0\) 时单调减少。函数在该邻域内并非单调减少。
\end{solutioncontent}

\mistakeknowledge{
\begin{itemize}
  \item \textbf{高阶导数极值判别法}：若 \(f(x)\) 在 \(x_0\) 的前 \(n-1\) 阶导数均为 0，但 \(f^{(n)}(x_0) \neq 0\)：当 \(n\) 为偶数时取极值（\(<0\) 为极大，\(>0\) 为极小）；当 \(n\) 为奇数时为拐点。
  \item \textbf{泰勒公式与极值/单调性判断}：判断增量 \(\Delta y\) 或导数符号时，可由泰勒展开中首个非零的高阶导数项直接决定。
\end{itemize}}

\end{mistake}
